\chapter{Arquitectura}

Este capítulo recoge por ahora el stack. Los componentes del orquestador están pendientes de implementar.

\section{Stack}

\begin{tabular}{ll}
  \textbf{Pieza} & \textbf{Elección} \\
  Lenguaje & Python 3.14 (D-13, D-124) \\
  Entorno y dependencias & \texttt{uv} (D-12) \\
  Protocolo & A2A 1.0 con \texttt{a2a-sdk} 1.3.0 (D-17) \\
  Almacén del estado & PostgreSQL 18.6 (D-69, D-125) \\
  Controlador de PostgreSQL & \texttt{psycopg} 3.3.8 (D-125) \\
  Tests & \texttt{pytest} (D-14) \\
  Lint y formato & \texttt{ruff} (D-15) \\
\end{tabular}

LangGraph resuelve con Python 3.14 (versión 1.2.20), pero aún no es una dependencia del proyecto. Se añade cuando se implementen los agentes.

\section{Comprobaciones de la base de datos}

Los tests de \texttt{tests/test\_setup.py} usan PostgreSQL real, sin simulaciones. Comprueban lo que el diseño pide a la base de datos:

\begin{itemize}
  \item Una transacción entre dos tablas se deshace entera si algo falla (D-63, D-69).
  \item Una escritura condicional por versión se aplica una sola vez: la segunda con la misma versión no afecta a ninguna fila (D-63).
  \item El reloj del servidor se puede leer y da la hora con zona horaria (D-69).
\end{itemize}

\section{Tasks de las peticiones del comprador}

D-97 exige escribir cada Task en la misma transacción que las sesiones. Los almacenes de \texttt{a2a-sdk} 1.3.0 no lo permiten, porque cada escritura abre y confirma su propia transacción. Por eso el orquestador guardará las Tasks en una tabla propia, con su propia implementación de \texttt{TaskStore} (D-126). Está pendiente de implementar.
